%% Causal Automata — arXiv / Zenodo / OSF preprint template
%% Style: kourgeorge/arxiv-style (MIT) — see paper/README.md
%% Edit the TODO blocks and compile with: make -C paper

\documentclass{article}

% --- Encoding / fonts / layout (arxiv.sty sets Times + geometry) ---
\usepackage[utf8]{inputenc}
\usepackage[T1]{fontenc}
\usepackage{arxiv}

\usepackage{amsmath,amssymb,amsfonts}
\usepackage{booktabs}
\usepackage{graphicx}
\usepackage{hyperref}
\usepackage{url}
\usepackage{natbib}
\usepackage{xcolor}
\usepackage{microtype}

% Figures live in figures/ relative to this file
\graphicspath{{figures/}}

% ---------------------------------------------------------------------------
% Metadata — fill these in first
% ---------------------------------------------------------------------------
\title{Causal Automata: GPU-Native Causal Discovery\\
via Synchronous Cellular Automata in JAX}

\author{
  Masakazu Yamagiwa\\
  Independent Researcher\\
  \texttt{myamagiwa@gmail.com}
  % \And
  % Coauthor Name\\
  % Affiliation\\
  % \texttt{email@example.com}
}

\date{\today}

% Optional header line (arxiv.sty)
\renewcommand{\headeright}{A Preprint}
\renewcommand{\undertitle}{A Preprint}
\renewcommand{\shorttitle}{Causal Automata}

\begin{document}
\maketitle

\begin{abstract}
%% TODO: 150--250 words. State the problem, method, and one quantitative claim.
We present \emph{causal-automata}, a synchronous cellular-automaton approach to
structure learning that recovers directed acyclic graphs (DAGs) and their Markov
equivalence classes (CPDAGs) from observational tabular data.
Candidate edges evolve under parallel birth and progressive
conditional-independence death, while a continuous node-elevation field
(V-well + gravity) orients the graph downhill.
Every update runs on a JAX / NumPy grid and compiles into a single XLA
\texttt{jax.lax.scan} loop on GPU/TPU.
On Benchpress ER and scale-free DAGs (\(p\in\{20,50,80,100\}\), \(n=1000\)),
causal-automata matches NOTEARS-level pattern F1 at roughly
\(40\times\) lower wall-clock time on an NVIDIA A10.
%% Soften or replace the claim above once you freeze the final numbers.
\end{abstract}

\keywords{causal discovery \and cellular automata \and DAG learning
  \and JAX \and Markov equivalence \and GPU}

% ===========================================================================
\section{Introduction}
\label{sec:intro}
% ===========================================================================

%% TODO: Motivation → gap → contribution bullets → roadmap.

Observational data yields correlations, not mechanisms.
Causal discovery aims to recover a DAG \(G=(V,E)\) whose edges encode directed
influence among measured variables~\citep{spirtes2000,pearl2009}.
Constraint-based methods such as PC~\citep{spirtes2000} and continuous
optimization methods such as NOTEARS~\citep{zheng2018} and
GOLEM~\citep{ng2020} dominate practice, but either scale poorly with
conditioning-set search or rely on sequential solvers that under-use modern
accelerators.

We treat adjacency discovery as a \emph{synchronous cellular automaton} on the
\(p\times p\) edge grid.
Local rules --- birth, death, V-well explaining-away, and elevation gravity ---
update every candidate edge in parallel, which maps naturally onto JAX/XLA.

\paragraph{Contributions.}
\begin{itemize}
  \item A physics-inspired CA formulation of causal discovery with explicit
        birth / death / V-well / gravity steps
        (Section~\ref{sec:method}).
  \item A JAX implementation that compiles the full run into one
        \texttt{jax.lax.scan} with static \texttt{ticks} and \(k\)
        (Section~\ref{sec:impl}).
  \item Empirical comparisons against PC, NOTEARS, and GOLEM on Benchpress
        and an in-repo gCastle harness (Section~\ref{sec:experiments}).
\end{itemize}

% ===========================================================================
\section{Background}
\label{sec:background}
% ===========================================================================

%% TODO: Keep short; point to related work for depth.

\paragraph{DAGs and CPDAG equivalence.}
A DAG encodes d-separation; under faithfulness, conditional independences in
the observational distribution identify the Markov equivalence class, often
represented as a CPDAG~\citep{spirtes2000}.

\paragraph{Conditional independence and colliders.}
Partial correlation (Fisher-\(z\)) tests approximate conditional independence
for linear-Gaussian data.
Colliders (V-structures) create dependence upon conditioning and supply an
orientation signal; we encode that signal as a downward force in elevation
space (the V-well).

\paragraph{Cellular automata and message passing.}
Pearl-style belief propagation aggregates local \(\pi\)/\(\lambda\)
messages~\citep{pearl1988}.
Cellular automata share the gather--aggregate--update motif; here the lattice
is the adjacency matrix itself rather than a spatial grid.

% ===========================================================================
\section{Method}
\label{sec:method}
% ===========================================================================

%% TODO: Replace sketch with your preferred notation from docs/theory.rst
%% and docs/physics.rst.

\subsection{State}
At tick \(t\) the automaton maintains:
\begin{itemize}
  \item binary (or fuzzy) adjacency \(A_t\in\{0,1\}^{p\times p}\),
  \item node elevations \(E_t\in[0,1]^p\),
  \item a fixed empirical covariance \(\mathrm{Cov}\) from the data matrix
        \(X\in\mathbb{R}^{n\times p}\).
\end{itemize}
Initialization is edgeless, \(A_0=0\), with a simple elevation schedule
(e.g.\ linspace) or learned warm-start.

\subsection{One tick}
Each tick applies, in order:
\begin{enumerate}
  \item \textbf{Gravity} --- project elevations so parents sit above children
        along the undirected skeleton \(U=\max(A,A^\top)\); optional multihop
        pulls toward two-hop means from \(A A\).
  \item \textbf{V-well} --- sink elevations at statistical colliders (with
        parent lift and Sepset screening against backdoor forks).
  \item \textbf{Death} --- prune edges that become conditionally independent
        given endpoint-local separator pools of size \(s=1..k\).
  \item \textbf{Birth} --- add edges that pass Sepset gates and a marginal
        correlation guard; orientation uses the elevation field.
\end{enumerate}
Optional fuzzy mode keeps continuous edge weights in \([0,1]\) and binarizes
only after the final tick.

\subsection{Output}
After \(T\) ticks we return the binarized DAG \(A_T\), elevations \(E_T\), and
optionally the CPDAG of the Markov equivalence class.

% ===========================================================================
\section{Implementation}
\label{sec:impl}
% ===========================================================================

%% TODO: Point to GitHub tag / version you freeze for the preprint.

The reference implementation is the open-source package
\texttt{causal-automata}~\citep{yamagiwa2026software}
(\url{https://github.com/kazuyamagiwa/causal-automata}).
NumPy (\texttt{bayesian.step}) and JAX (\texttt{jax\_core.step} /
\texttt{run}) share the same physics.
The JAX path marks \texttt{ticks} and \(k\) as static arguments so scan length
and partial-correlation block sizes are compile-time constants.

% ===========================================================================
\section{Experiments}
\label{sec:experiments}
% ===========================================================================

%% TODO: Replace with your frozen Benchpress / gCastle tables and figures.
%% Drop PNGs into figures/ and \includegraphics them.

\subsection{Setup}
We evaluate on Benchpress~\citep{benchpress} Erdős--Rényi and scale-free DAGs
with \(p\in\{20,50,80,100\}\), mean degree \(d=4\), \(n=1000\), and seeds
\(1..10\) (80 cells), on an NVIDIA A10.
Baselines: PC, NOTEARS, GOLEM (gCastle / Benchpress configurations --- document
exact versions here).

\subsection{Results}
Table~\ref{tab:benchpress} sketches the headline numbers from the software
README (v0.3.6); replace with the table you want frozen in the preprint.

\begin{table}[t]
  \centering
  \caption{Benchpress summary (placeholder --- verify before submission).}
  \label{tab:benchpress}
  \begin{tabular}{lrrrr}
    \toprule
    Algorithm & time (s) & F1 skel. & F1 pattern & SHD pattern \\
    \midrule
    Causal Automata & \(0.76\pm0.32\) & \(0.822\pm0.088\) & \(0.754\pm0.121\) & \(50.1\pm40.3\) \\
    NOTEARS         & \(31.4\pm29.5\) & \(0.896\pm0.056\) & \(0.768\pm0.092\) & \(39.0\pm18.9\) \\
    PC              & \(80.4\pm342\)  & \(0.888\pm0.061\) & \(0.720\pm0.085\) & \(50.4\pm27.5\) \\
    GOLEM           & \(18.5\pm11.2\) & \(0.679\pm0.086\) & \(0.469\pm0.084\) & \(95.2\pm55.5\) \\
    \bottomrule
  \end{tabular}
\end{table}

%% Example figure slot (uncomment when you add a PNG):
% \begin{figure}[t]
%   \centering
%   \includegraphics[width=0.95\linewidth]{elevation_convergence}
%   \caption{From empty grid to recovered DAG (ground-truth ghost in grey).}
%   \label{fig:hero}
% \end{figure}

% ===========================================================================
\section{Discussion and limitations}
\label{sec:discussion}
% ===========================================================================

%% TODO: Faithfulness, linearity, continuous data, hyperparameter \(k\), etc.

Assumptions include (approximate) faithfulness and, for Fisher-\(z\) screens,
linear-Gaussian structure.
Elevation orientation can err under near-ties; fuzzy mode and multihop gravity
trade compute for accuracy.
We leave discrete / mixed-type data and interventional settings to future work.

% ===========================================================================
\section{Conclusion}
\label{sec:conclusion}
% ===========================================================================

Causal automata cast structure learning as synchronous local updates that map
onto GPU scan loops.
The accompanying JAX package makes the method practical; this preprint is the
narrative companion.
%% TODO: one-sentence takeaway you want cited.

% ===========================================================================
\section*{Acknowledgments}
% ===========================================================================

%% TODO: funding, compute credits, feedback.
We thank early users and benchmark maintainers of Benchpress and gCastle.

% ===========================================================================
\section*{Data and code availability}
% ===========================================================================

Source code, docs, and regeneration scripts:
\url{https://github.com/kazuyamagiwa/causal-automata}.
Package on PyPI: \texttt{causal-automata}.
%% TODO: Zenodo DOI for the frozen tarball / this PDF once deposited.

\bibliographystyle{plainnat}
\bibliography{references}

\appendix
\section{Notation}
\label{app:notation}
\begin{center}
\begin{tabular}{ll}
  \toprule
  Symbol & Meaning \\
  \midrule
  \(p\) & number of variables \\
  \(n\) & number of samples \\
  \(A\) & adjacency matrix \\
  \(E\) & node elevations \\
  \(k\) & max conditioning-set size \\
  \(T\) & number of CA ticks \\
  \bottomrule
\end{tabular}
\end{center}

\section{Hyperparameters}
\label{app:hparams}
%% TODO: list ticks, k, thresholds, seeds, hardware.
Default estimator: \texttt{CausalAutomaton(ticks=100, k=3)}.

\end{document}
